\documentclass[a4paper, serif, 10pt]{moderncv}

%% moderncv
% color and style
\moderncvcolor{black}
\moderncvstyle{banking}

%% page layout/margins
\usepackage[top=2.5cm, bottom=2.5cm, left=1.5cm, right=1.5cm]{geometry}

\nopagenumbers{}

%% Babel config for Italian language
% ref:
% - https://latex3.github.io/babel/guides/locale-italian.html
\usepackage[italian]{babel}

%% fontspec
\usepackage{fontspec}

\IfFontExistsTF{Linux Libertine}{
  \setmainfont[Ligatures={TeX, Common}, Numbers=OldStyle]{Linux Libertine}
}{}
\IfFontExistsTF{Linux Libertine O}{
  \setmainfont[Ligatures={TeX, Common}, Numbers=OldStyle]{Linux Libertine O}
}{}

%% CTeX
% CJK support, used to show CJK characters in the resume
%
% - heading=false: disable ctex heading
% - punct=kaiming: use kaiming punctuations styles for CJK
% - scheme=plain: use plain scheme, do not override \normalsize font size
% - space=auto: space settings for CJK characters
%
% CJK fonts are resolved per locale language, e.g. Noto Serif/Sans CJK JP
% for Japanese resumes, with fallback to the next region if not installed
%
% ref:
% - http://ctan.mirrorcatalogs.com/language/chinese/ctex/ctex.pdf
\usepackage[UTF8, heading=false, punct=kaiming, scheme=plain, space=auto]{ctex}

\IfFontExistsTF{Noto Serif CJK SC}{
  \setCJKmainfont{Noto Serif CJK SC}
}{}
\IfFontExistsTF{Noto Sans CJK SC}{
  \setCJKsansfont{Noto Sans CJK SC}
}{}

%% line spacing
\usepackage{setspace}
\setstretch{1.125}

%% URL styling - use normal text instead of monospace
\urlstyle{same}

%% Patch \httplink and \httpslink to handle full URLs
% The original moderncv macros always prepend http:// or https:// to URLs.
% This causes issues when the URL already has a protocol (e.g., https://example.com)
% resulting in malformed URLs like https://https://example.com.
% This patch checks if the URL already contains "://" and uses it directly if so.
% Note: We use \str_set:Nx (x-expansion) to expand macros like \@homepage first.
\ExplSyntaxOn
\str_new:N \g_yamlresume_url_str

\RenewDocumentCommand{\httpslink}{O{}m}{
  \str_set:Nx \g_yamlresume_url_str {#2}
  \str_if_in:NnTF \g_yamlresume_url_str {://}
    { \url{#2} }
    { \url{https://#2} }
}

\RenewDocumentCommand{\httplink}{O{}m}{
  \str_set:Nx \g_yamlresume_url_str {#2}
  \str_if_in:NnTF \g_yamlresume_url_str {://}
    { \url{#2} }
    { \url{http://#2} }
}
\ExplSyntaxOff

\name{Luca Bianchi}{}
\title{Data Engineer | Piattaforme Dati e Analytics Cloud}
\phone[mobile]{+39 346 123 4567}
\email{luca.bianchi.data@gmail.com}
\homepage{https://lucabianchi.dev}

\address{Via Tortona 31, Milano, Lombardia, Italia, 20144}{}{}

\extrainfo{{\small \faLinkedin}\ \href{https://www.linkedin.com/in/luca-bianchi-data}{@luca-bianchi-data} {} {} {} • {} {} {} 
{\small \faGithub}\ \href{https://github.com/lucabianchi-data}{@lucabianchi-data}}

\begin{document}

\maketitle

\section{Informazioni personali}

\cvline{}{Data Engineer con oltre 9 anni di esperienza nella progettazione e nello sviluppo di piattaforme dati scalabili su cloud.
Specializzato in pipeline ETL/ELT, modellazione dei dati e orchestrazione con Apache Airflow e dbt.
Esperienza consolidata con AWS, Snowflake e BigQuery, maturata in contesti Agile e team interfunzionali.
Orientato alla qualità del dato, all'automazione e al monitoraggio delle pipeline.
Appassionato di open source e della diffusione della cultura data-driven.}
\section{Formazione}

\cventry{set 2014–lug 2016}
        {Laurea magistrale, Ingegneria Informatica, Punteggio: 110/110 e lode}
        {Politecnico di Milano}
        {\url{https://www.polimi.it}}
        {}
        {\begin{itemize}
\item Tesi magistrale su architetture dati distribuite per l'analisi di grandi volumi di dati
\item Approfondimento di basi di dati relazionali e NoSQL, modellazione dimensionale e processi ETL
\item Partecipazione a progetti di gruppo con metodologie Agile e presentazioni tecniche
\end{itemize}
\textbf{Corsi}: Basi di Dati e Sistemi Informativi, Ingegneria del Software, Sistemi Distribuiti, Machine Learning, Cloud Computing, Architetture dei Sistemi di Elaborazione, Sicurezza dei Sistemi Informatici}

\cventry{set 2011–lug 2014}
        {Laurea triennale, Ingegneria Informatica, Punteggio: 108/110}
        {Università degli Studi di Milano}
        {\url{https://www.unimi.it}}
        {}
        {\begin{itemize}
\item Solida preparazione in programmazione, algoritmi e basi di dati
\item Esperienza pratica con Java, Python e SQL in progetti universitari
\item Sviluppo di competenze di analisi e problem solving
\end{itemize}
\textbf{Corsi}: Fondamenti di Programmazione, Algoritmi e Strutture Dati, Basi di Dati, Reti di Calcolatori, Sistemi Operativi, Analisi Matematica}
\section{Esperienza professionale}

\cventry{gen 2022–Presente}
        {Data Engineer Senior}
        {FintechLab Italia}
        {\url{https://www.fintechlabitalia.it}}
        {}
        {\begin{itemize}
\item Progettazione e implementazione di una piattaforma dati cloud-native su AWS con Snowflake e dbt
\item Orchestrazione di oltre 120 pipeline giornaliere con Apache Airflow, riducendo i tempi di elaborazione del 35\%
\item Definizione di standard di qualità del dato e osservabilità con Great Expectations e monitoraggio su Grafana
\item Collaborazione con team prodotto e compliance per la modellazione di dati finanziari e reportistica regolamentare
\item Mentoring di due data engineer junior e revisione delle pull request
\end{itemize}
\textbf{Parole chiave}: Data Platform, Apache Airflow, dbt, Snowflake, AWS, Data Quality}

\cventry{feb 2019–dic 2021}
        {Data Engineer}
        {Retail Analytics S.p.A.}
        {\url{https://www.retailanalytics.it}}
        {}
        {\begin{itemize}
\item Sviluppo di pipeline ETL/ELT per l'integrazione di dati da CRM, ERP e piattaforme e-commerce
\item Modellazione di un data warehouse su Google BigQuery con viste materializzate e partizionamento
\item Automazione dei flussi con Airflow e Python, con test automatici e deploy tramite CI/CD
\item Ottimizzazione delle query SQL e riduzione dei costi di elaborazione del 20\%
\item Realizzazione di dashboard self-service in Looker per i team commerciale e marketing
\end{itemize}
\textbf{Parole chiave}: ETL/ELT, Google BigQuery, Python, Apache Airflow, Looker, Data Warehouse}

\cventry{set 2016–gen 2019}
        {Business Intelligence Analyst}
        {Consulenza Dati Milano}
        {\url{https://www.consulenzadatimilano.it}}
        {}
        {\begin{itemize}
\item Raccolta dei requisiti e sviluppo di report e dashboard in Power BI per clienti del settore manifatturiero
\item Scrittura di query SQL complesse su SQL Server e PostgreSQL per analisi operative e direzionali
\item Supporto alla definizione di KPI e modelli dati con i referenti business
\item Documentazione dei processi di reporting e formazione degli utenti finali
\end{itemize}
\textbf{Parole chiave}: Power BI, SQL Server, PostgreSQL, KPI, Data Visualization}
\section{Lingue}

\cvline{Italiano}{Madrelingua o bilingue \hfill \textbf{Parole chiave}: Madrelingua}
\cvline{Inglese}{Competenza professionale completa \hfill \textbf{Parole chiave}: IELTS 7.5, C1 Business English}
\cvline{Francese}{Competenza professionale limitata \hfill \textbf{Parole chiave}: DELF B1}
\section{Competenze}

\cvline{Data Engineering}{Esperto \hfill \textbf{Parole chiave}: Python, SQL, Apache Spark, PySpark, Apache Airflow, dbt, Kafka, ETL/ELT}
\cvline{Cloud \& Data Warehouse}{Avanzato \hfill \textbf{Parole chiave}: AWS, Snowflake, Google BigQuery, Azure, PostgreSQL, Docker, Kubernetes, Terraform}
\cvline{Analytics \& BI}{Avanzato \hfill \textbf{Parole chiave}: Power BI, Tableau, Looker, Data Modeling, Data Quality, Data Governance}
\cvline{Linguaggi e Strumenti}{Avanzato \hfill \textbf{Parole chiave}: Scala, Java, Bash, Git, Jira, Confluence, Agile/Scrum}
\section{Premi}

\cventry{nov 2023}
        {Data Summit Italia}
        {Premio Data Innovation}
        {}
        {}
        {Riconoscimento per il progetto di osservabilità delle pipeline dati che ha ridotto del 40\% il tempo medio di risoluzione degli incidenti.}

\cventry{dic 2020}
        {Retail Analytics S.p.A.}
        {Miglior Progetto Dati dell'Anno}
        {}
        {}
        {Premio interno per la realizzazione del data warehouse su BigQuery e delle dashboard self-service per il team commerciale.}
\section{Certificazioni}

\cventry{mar 2024}
        {Amazon Web Services}
        {AWS Certified Data Engineer - Associate}
        {\url{https://aws.amazon.com/certification/certified-data-engineer-associate/}}
        {}
        {}

\cventry{giu 2022}
        {Google Cloud}
        {Google Professional Data Engineer}
        {\url{https://cloud.google.com/learn/certification/data-engineer}}
        {}
        {}

\cventry{set 2023}
        {dbt Labs}
        {dbt Analytics Engineering Certification}
        {\url{https://www.getdbt.com/certifications/analytics-engineer/}}
        {}
        {}
\section{Pubblicazioni}

\cventry{apr 2023}
        {Pipeline dati affidabili: pattern di osservabilità per team analytics}
        {Rivista Italiana di Data Engineering}
        {\url{https://www.ridata.it/articoli/osservabilita-pipeline-dati}}
        {}
        {\begin{itemize}
\item Articolo tecnico su metriche di qualità, monitoraggio e gestione degli incidenti nelle pipeline dati
\item Condivisione di pattern applicabili con Airflow, dbt e Great Expectations
\item Casi pratici tratti da contesti fintech e retail
\end{itemize}}
\section{Progetti}

\cventry{gen 2023–giu 2023}
        {Framework open source per il monitoraggio della qualità dei dati nelle pipeline Airflow.}
        {Data Quality Radar}
        {\url{https://github.com/lucabianchi-data/data-quality-radar}}
        {}
        {\begin{itemize}
\item Progettazione di controlli di qualità riutilizzabili per tabelle Snowflake e BigQuery
\item Integrazione con Airflow per l'esecuzione schedulata e la notifica degli alert
\item Documentazione e rilascio su PyPI con oltre 300 download mensili
\end{itemize}
\textbf{Parole chiave}: Data Quality, Great Expectations, Apache Airflow, Python}

\cventry{mar 2021–nov 2022}
        {Piattaforma interna per l'accesso governato ai dati e la creazione di dashboard self-service.}
        {Self-Service Analytics Hub}
        {\url{https://github.com/lucabianchi-data/analytics-hub}}
        {}
        {\begin{itemize}
\item Modellazione di dataset certificati in dbt e Snowflake per i team business
\item Definizione di ruoli e permessi per l'accesso ai dati sensibili
\item Formazione di oltre 60 utenti su Looker e sulle best practice di analytics
\end{itemize}
\textbf{Parole chiave}: Analytics, Looker, Snowflake, dbt}
\section{Interessi}

\cvline{Open Source}{Apache Airflow, dbt, Great Expectations}
\cvline{Sport}{Ciclismo, Nuoto, Trekking}
\cvline{Fotografia}{Paesaggio, Street photography}
\section{Volontariato}

\cventry{set 2021–Presente}
        {Mentore Data Engineering}
        {Women in Tech Italia}
        {\url{https://www.womenintechitalia.it}}
        {}
        {\begin{itemize}
\item Mentoring per donne che intraprendono percorsi di carriera nel data engineering
\item Supporto nella preparazione di colloqui tecnici e nella definizione di piani di crescita
\item Partecipazione a incontri di community e workshop su SQL e pipeline dati
\end{itemize}}
    

\cventry{gen 2020–dic 2023}
        {Volontario e speaker}
        {Data Science Milano}
        {\url{https://www.datasciencemilano.it}}
        {}
        {\begin{itemize}
\item Organizzazione di meetup mensili su data engineering, analytics e machine learning
\item Interventi su orchestrazione di pipeline e data quality
\item Coordinamento di volontari e relazioni con speaker e sponsor
\end{itemize}}
    
\end{document}